% !LW recipe = make-deck
% Level 3T (proofs track, TPIL): difference, proved. Theory: difference,
% intersection, law 4, and a fact with a classical step. Lean: ¬, and
% classical logic.
% Build: make (rules in ../common/deck.mk).
\DocumentMetadata{
  lang        = en-GB,
  pdfstandard = ua-2,
  tagging     = on,
}
\documentclass{ltx-talk}

\input{preamble}

\title{Level 3T: Difference, Proved}
\author{The Proofs and Programs Club}
\institute{Royal Holloway, University of London}

\begin{document}

\leveltitle{3T}{Difference, Proved: ``Not'', and Classical Logic}
  {Proofs track: follows TPIL}{After \levellink{2t}}
  {Next: \levellink{4t}}

\levelframe{3T}{Level 3T at a Glance}{%
  difference and intersection, on tables as properties of rows.}{%
  law 4, and that $R \cap S = R - (R - S)$: one half of it needs
  classical logic.}{%
  prove and use ``not'': $\neg A$ means ``if $A$, then false'' &
    \rds{tpil-proplogic} \\
  reason classically: prove $A$ by refuting $\neg A$ &
    \rds{tpil-classical} \\
}

% ---------------------------------------------------------------- theory
\theorypart{Difference and intersection}

\begin{frame}
  \frametitle{Difference Removes the Rows of Another Table}
  \defn{Definition 3T.1 (Difference, intersection)}{For tables $R$ and
    $S$ of the same rows: $R - S = \{t \in R \mid t \notin S\}$ and
    $R \cap S = \{t \in R \mid t \in S\}$.}
  Which sessions had nobody? With one-column tables:
  \[
    \pi_1[\rel{Session}] - \pi_2[\rel{Attended}] = \{1, 2, 3, 4\} - \{1,
    2, 3\} = \{4\}.
  \]
  \defn{Law 4}{$\sel{p}(R - S) = \sel{p}(R) - S$.}
  \emph{Proof.} $t \in \sel{p}(R - S) \iff t \in R$, $t \notin S$ and
  $p(t)$ $\iff t \in \sel{p}(R)$ and $t \notin S$ $\iff t \in \sel{p}(R)
  - S$. \proofend
  \slideref{\lnk{ahv5}}
  \theorytag{3T}
\end{frame}

\begin{frame}
  \frametitle{Intersection Comes from Difference, Classically}
  \defn{Definition 3T.2 (Not)}{``Not $A$'' means: if $A$, then a
    contradiction. So $t \notin S$ means: if $t \in S$, then a
    contradiction.}
  \defn{Fact 3T.3}{$R \cap S = R - (R - S)$.}
  \emph{Proof.} $t \in R - (R - S)$ means: $t \in R$, and not ($t \in R$
  and $t \notin S$).
  \begin{itemize}
    \item If $t \in R$ and $t \in S$: then not ($t \in R$ and
      $t \notin S$), since $t \notin S$ would contradict $t \in S$.
    \item If $t \in R$ and not ($t \in R$ and $t \notin S$): then
      $t \notin S$ is impossible, so \textbf{classically} $t \in S$.
      \proofend
  \end{itemize}
  The last step uses: if not-$A$ is impossible, then $A$. It is not
  provable from Definition 3T.2 alone.
  \theorytag{3T}
\end{frame}

% ---------------------------------------------------------------- Lean
\leanpart{Difference, and a classical proof}

\begin{frame*}
  \frametitle{Difference and Intersection Are Formulas Too}
  \leanblock{diff.lean}
  \begin{itemize}
    \item Put this after your code of Level 2T: \code{Table} and
      \code{select} come from there.
    \item \code{¬ A} is \code{A → False}: a function that turns a proof
      of \code{A} into a proof of \code{False}, which has no proof.
  \end{itemize}
  \slideref{\lnk{tpil-proplogic}}
  \leantag{3T}
\end{frame*}

\begin{frame*}
  \frametitle{Half of Fact 3T.3, with No Classical Logic}
  \leanblock{first.lean}
  \begin{itemize}
    \item \code{fun ⟨hr, hs⟩ => \ldots} takes the pair apart: \code{hr}
      proves \code{R t}, and \code{hs} proves \code{S t}.
    \item The second half of the answer proves a \code{¬}: it receives a
      proof of \code{diff R S t}, whose \code{hns : ¬ S t} meets
      \code{hs}.
  \end{itemize}
  \slideref{\lnk{tpil-proplogic}}
  \leantag{3T}
\end{frame*}

\begin{frame*}
  \frametitle{Your Turn: Law 4, and the Classical Half}
  \leanblock{exercises.lean}
  \begin{itemize}
    \item \ex{3T.1} Law 4. \ex{3T.2} The classical half, with
      \code{Classical.byContradiction : (¬ A → False) → A}.
  \end{itemize}
  \slideref{\lnk{tpil-classical}}
  \leantag{3T}
\end{frame*}

\end{document}
